\documentclass[12pt]{article}

\usepackage[T1]{fontenc}
\usepackage[utf8]{inputenc}
\usepackage[margin=1in]{geometry}
\usepackage{setspace}
\usepackage{parskip}
\usepackage{titlesec}
\usepackage[hidelinks]{hyperref}

\titleformat{\section}{\normalfont\bfseries\large}{}{0pt}{}
\titlespacing*{\section}{0pt}{1.4em}{0.6em}

\setstretch{1.15}

\begin{document}

\begin{center}
	{\Large\bfseries Primary Source Worksheet \#1}\\[0.6em]
	RELI 252 \quad\textbullet\quad Dr.\ Clarke\\[0.3em]
	Greyson Knippel
\end{center}

\vspace{1em}
\hrule
\vspace{1.5em}

\noindent\textbf{Primary Source Title:} ``An Address to the Negroes in the State of New York''

\vspace{0.5em}

\noindent\textbf{Primary Source Author:} Jupiter Hammon (delivered 1786, published 1787)

\section*{What is the author's intent with this piece?}

Hammon wants his audience to change their priorities. He wants them to treat salvation as their main goal, put freedom as a secondary goal, and to accept obedience as some sort of a religious obligation rather than to have agency over their own lives. He hopes listeners are scared of damnation, convinced their suffering is temporary, and willing to behave in a way that won't ruin the chances for gradual emancipation that he actually does support.

He also sets himself up as someone worth listening to. He calls his speech ``\ldots the last, and I may say dying advice, of an old man,'' which makes his criticism sound like care instead of judgment. He also points out exactly why his audience should trust him. He says they ``\ldots will be more likely to listen to what is said, when you know it comes from a negro, one of your own nation and colour, and therefore can have no interest in deceiving you.'' He is using his own race as proof that he is honest, which is something a white preacher could never do.

Most of his argument is built on fear, not logic. He spends a long time on Judgment Day, when everything done in secret ``shall be declared on the house top.'' He describes a hell his listeners ``laugh at'' now but will be stuck in forever. Then he makes slavery seem small next to that by asking, ``What is forty, fifty, or sixty years, when compared to eternity?'' He never says slavery is fair. He simply just makes the question feel unimportant compared to eternity.

Hammon is also partially writing for white readers. He mentions that white people read his earlier work, and he ends by telling free Blacks that bad behavior ``prevent[s] our being free.'' That part is clearly aimed at the white people deciding if emancipation is a good idea.

\section*{Compare and contrast with a previously assigned primary source}

Hammon's address and Phillis Wheatley's ``On Being Brought from Africa to America'' use almost the same religious language. Both writers were enslaved, both write inside white Protestant Christianity, and both say that being exposed to Christianity saved them from something worse. Wheatley writes that ```Twas mercy brought me from my Pagan land.'' Hammon tells his listeners to ``bless God for ever, that you have been brought into a land where you have heard the gospel, though you have been slaves.'' Neither one treats African religion as something worth keeping, which is surprisingly different from how back when Equiano writes about his Igbo background.

The big difference though is who they are talking to. Wheatley is talking to white Christians and calling them out. She writes, ``Remember, Christians, Negros, black as Cain, / May be refin'd, and join th' angelic train.'' She uses the idea that anyone can be saved to argue that Black people are spiritually equal. Hammon is talking to Black listeners and correcting them about stealing, swearing, and disobedience while white readers watch. Both use the same religion, but Wheatley uses it to argue for equality and Hammon uses it to argue for obedience.

Their situations explain a lot of this. Hammon was around seventy, enslaved his whole life, and says he does ``not wish to be free.'' He also doubts that older slaves could ``take care of ourselves.'' Wheatley was freed young and was famous on both sides of the Atlantic. She could picture Black freedom because she already had hers. Hammon is also writing in 1786, right after the Revolution, and he notices all the liberty talk around him. He points out ``how many lives have been lost to defend their liberty,'' then redirects that energy toward heaven instead of emancipation.

Together, these show that early African American Christianity was not one single position. The same Bible was used to argue for patience and for equality, depending on who was speaking and who was listening. Lots of people had lots of different experiences, and those experiences massively change your views of the world, even if you come from the same religion and culture.

\end{document}